\cotefile{157-3}
\scanpage{342}
\mspage{41}

de l’âme est sans commencement et sans fin ! Cela signifie, par contre, libérer peu ou prou ce désir des vibrations parasites issues de l’égo qui le brouillent. Parmi celles-ci, il y a surtout celles qui émanent de la redoutable figure de \textit{Censeur}, gardien des répressions qui s’évertuent à cacher ce désir et à le désavouer. Avec chaque libération partielle, instaurant notre autonomie par rapport à l’une ou l’autre des figures archétypes, le Censeur s’amenuise, il perd de son autorité et de sa prestance. Et quand nous sommes libres enfin et de nos Parents, et de nos Conjoints, le Censeur n’est plus que l’ombre de lui-même ! Une fois la guerre des générations et la guerre du sexe terminée, la mort du Censeur ne saurait plus tarder...

Parmi les quatre pulsions amoureuses archétypes, nous attirant en le corps de la Mère, en le corps de l’Epouse, en le corps du Père et en celui de l’Epoux, deux sont incestueuses, et deux sont homosexuelles - et pourtant elles sont ! Et elles sont éternellement créatrices et éternellement belles. Une seule parmi les quatre (celle qui porte l’homme pubère vers l’Epouse, la femme nubile vers l’Epoux) est censée être recevable, ou ne serait-ce qu’\textit{exister} selon le Plan de Dieu, ou selon les lois de Nature. Mais il n’en est pas une seule qui ne soit en Dieu de toute éternité, et qu’Il n’ait insufflée dans l’âme dès sa création, pour animer son élan vers ses destinées éternelles.

\textit{La grande tâche de l’âme au temps de son errance, celle qui précède et prépare toutes les autres, c’est de se libérer de sa dépendance par rapport à ses Amants et par rapport à ses Amantes originels, et ce faisant, d’apprendre à aimer.} Et aimer dans la plénitude, c’est aussi prendre connaissance du désir éternel qui nous porte vers chacun d’eux, et assumer ce désir dans toute sa force et dans toute sa beauté. Quand cette tâche-là est remplie, l’âme est prête enfin pour entrer dans les splendeurs éternelles.

\noindent 17. Assumer le désir, ce n’est pas nécessairement l’assouvir

Assumer un désir charnel, ce n’est ni plus ni moins qu’en prendre connaissance au niveau conscient, et d’en sentir intensément la beauté. C’est aussi ``\textit{vivre}'' ce désir dans sa force et dans sa beauté originelle, le vivre comme il est vécu par Dieu en Lui-même, et comme il est vécu aussi par les couches profondes de l’âme, en certains moments privilégiés. Lorsque (comme il est presque toujours le cas) le désir se heurte à des interdits intériorisés par la psyché, \textit{assumer ce désir est un acte de vérité et de courage, un acte de fidélité}, dont la portée est d’autant plus grande que les interdits sont plus puissants. Dans la progression spirituelle de l’âme, ces actes-là ne jouent pas un rôle moins crucial que ceux qui tendent à la libérer de la vanité. Les uns tendent à instaurer une autonomie par rapport au ``\textit{Flatteur}'', les autres une autonomie non moins ardue par rapport au Censeur.

Il est important de réaliser qu’assumer, donc aussi ``vivre'' un désir, est un acte d’essence différente de l’acte d’assouvissement. Aucun des deux actes n’implique l’autre : nous pouvons vivre un désir, ``vibrer'' en lui dans une plénitude parfaite, sans jamais l’assouvir ``en actes'' ni même (si les circonstances sont adverses) souhaiter l’assouvir ainsi, charnellement. Inversement, même les désirs charnels ressentis comme parfaitement orthodoxes sont presque toujours assouvis sans être véritablement assumés. Il en est ainsi quand les comportements amoureux sont dominés par les poses conventionnelles (ou d’autres poses de nature plus personnelles), imposées par la sempiternelle immixition de l’égo dans les ébats amoureux.

% archive sheet number: 342
\scanpage{343}
\mspage{\unsure{9|1}}

\diagram{\begin{tikzcd}[ampersand replacement=\&,column sep=small,row sep=small]
{} \&[-.4em] {} \& {} \& b \& {} \\
\deleted{X}V={} \& \biggl(\unsure{0|o}\arrow[r,bend left=16,"\lambda"]\arrow[r,bend right=16,"\mu"'] \& a\arrow[ur,"\beta"]\arrow[dr,"\gamma"'] \& {} \& \biggr)\\
{} \& {} \& {} \& c
\end{tikzcd}}

avec
\[
\beta\lambda=\beta\mu\qquad\text{et}\qquad\gamma\lambda=\gamma\mu
\]
les deux seules flèches non identiques, en plus des \inserted{quatre} flèches marquées $\deleted{\unsure{B}}\beta,\deleted{\unsure{\beta}}\gamma,\deleted{\unsure{\beta}}\lambda,\mu$.

On trouve que \unsure{!|1}
\[
b\amalg_a c \overset{\sim}{\longrightarrow}e_X
\qquad\text{(somme amalgamée dans $\deleted{X}V^{\wedge}$)}
\]

Pourtant $\deleted{X}V$ n’est pas asphérique, ni même $1$-connexe : il \deleted{\unsure{a un|a le} \unsure{groupe|même}} \inserted{est} homotope \unsure{à|:} la sous-catégorie \unsure{ouverte|ouvert}
\diagram{\begin{tikzcd}[ampersand replacement=\&,column sep=small]
\deleted{X}V_0={} \& 0\arrow[r,bend left=18,"\lambda"]\arrow[r,bend right=18,"\mu"'] \& a,
\end{tikzcd}}
\deleted{\unsure{lorsq}} car on a une rétraction $\deleted{X}V\xrightarrow{r}\deleted{X}V_0$ (\unsure{envoie|envoyant} $b$ et $c$ sur $a$), et \deleted{$ir$} une flèche $ir\longrightarrow\mathrm{id}_{\deleted{X}V}$ (\unsure{et|où} $i:X_0\unsure{\hookrightarrow|\subset}X$ est l’inclusion). Donc $\deleted{X}V$ a le $=$ type d’homotopie que \uline{$\deleted{X}V_0$}, qui est un cercle.

Donc on a intérêt à prendre ce $V$, \inserted{\unsure{ou|où} $\widetilde V$ déduit de $V$ en rajoutant un élément final} \deleted{\unsure{comme}}
\diagram{\begin{tikzcd}[ampersand replacement=\&,column sep=small,row sep=small]
{} \& {} \& {} \& b\arrow[dr] \\
\widetilde V={} \& 0\arrow[r,bend left=18]\arrow[r,bend right=18] \& a\arrow[ur]\arrow[dr] \& {} \& e,\\
{} \& {} \& {} \& c\arrow[ur]
\end{tikzcd}}
comme $\unsure{W|W'}$, pour les contre-exemples page $6$ et $7$\deleted{\unsure{.}}

% archive sheet number: 343
\scanpage{344}
\mspage{44}

celle-ci est le plus souvent non moins cachée dans l’inconscient que celle-là, tout en s’exprimant de façon tout aussi flagrante dans vos actes et dans vos comportements.

\noindent 19. La violence de rétaliation et les préceptes de Jésus

Il me reste enfin à parler encore d’un type d’agressivité, ou de violence, qui n’est qu’indirectement liée aux pulsions d’antagonisme vis-à-vis des archétypes fondamentaux. Il s’agit de ce qu’on pourrait appeler la ``\textit{violence de réaction}''. C’est la violence qui naît dans la psyché en réaction à des agressions subies, réaction d’autant plus forte que ces agressions sont ressenties comme injustifiées (par exemple, comme excessives), voire même comme totalement ``gratuites''.

La violence de réaction prend souvent la forme d’une \textit{rétaliation}, s’exerçant à l’encontre de l’agresseur (réel ou présumé). Encore faut-il que l’agresseur soit connu et qu’il soit vulnérable. Il n’en est rien pour les ``\textit{agressions diffuses}'', c’est-à-dire pour celles qui ne sont pas le fait de telle personne en particulier, mais de tout un milieu ou de conditions de vie, reconnues par les couches profondes comme incompatibles avec le bien-être et l’épanouissement de la psyché. Les agressions subies tout au cours de l’enfance de la part du milieu parental sont le plus souvent de cet ordre. Un autre cas fréquent où une libération des tensions psychiques par une violence de retaliation est freinée, et souvent même réduite à des actes purement symboliques et de faible effet sur la psyché, est celui où l’agresseur est ressenti (à tort ou à raison) comme trop puissant pour qu’une agression contre lui soit pensable. (Si ce n’est tout au plus dans des occasions propices, quand il est ressenti comme étant en position de faiblesse, et quand l’âme croit pouvoir frapper sans risques...)

Mais alors même que l’agressivité de réaction accumulée par la psyché ne pourrait trouver exutoire et se libérer par des actes de rétaliation, elle trouve nécessairement une expression erratique, par ce qu’on peut appeler la ``\textit{violence de défoulement}''. C’est encore une forme de violence gratuite, frappant au gré des occasions partout où elle peut frapper avec un minimum de risques de rétaliation. Cette forme de violence prolifère surtout en périodes de crise, ou dans des sociétés en voie de décomposition. Elle fait partie des nombreux ``signes de la Fin'', qui marquent de plus en plus au fil des ans et des décennies, le ``\textit{Temps de la Fin}'' dans lequel nous sommes engagés depuis la première guerre mondiale, et qui va toucher à sa fin le 14 Octobre 1996.

Les mécanismes de la ``violence de réaction'', par quoi toute violence subie suscite une violence en retour, fait partie des \textit{conditionnements originels de l’âme}. La libération de la dépendance par rapport aux Figures archétypes, laquelle met fin à la ``violence automatique'' (ou de transfert) liée à ces figures, ne met pas fin pour autant à ce conditionnement plus tenace et plus impérieux encore que cette dépendance. Comme pour tout autre conditionnement originel, la ``vulnérabilité à l’agression'' ne prend fin que par un Acte de Dieu, entérinant les efforts faits par l’âme pour \unsure{le} dépasser. Le plus souvent, cet Acte de Dieu s’accomplit par l’Acte de libération. (Cet Acte met fin à chacun des conditionnements originels qui restent actifs dans l’âme, à la seule exception de la répression sexuelle.)

De cette vulnérabilité de l’âme à la violence, la plupart des messagers du Nouvel Age (et parmi eux la plupart des destinataires de la Lettre) ne seront pas libérés dès leur présente existence, mais seulement au terme

% archive sheet number: 344
\scanpage{345}
\mspage{\deleted{\unsure{9|10}}10}

2.6. Contre-exemples pour la totale asphéricité.

Je me suis finalement convaincu que
\diagram{\begin{tikzcd}[ampersand replacement=\&,row sep=small,column sep=small]
{} \& {} \& b_0\\
\unsure{\Upsilon|Y}=\biggl\{ \& a_0\arrow[ur,"\beta_0"]\arrow[dr,"\gamma_0"'] \& {} \& \biggr\}\\
{} \& {} \& c_0
\end{tikzcd}}
n’est \uline{pas} totalement asphérique (pour le localiseur $W_\infty$ -- dans le présent numéro, on supposera $W=W_\infty$), ou ce qui revient au $=$, qu’on peut avoir un diagramme de type $\Upsilon$
\diagram{\begin{tikzcd}[ampersand replacement=\&,row sep=small,column sep=small]
{} \& {} \& b\\
\text{(a)} \& a\arrow[ur,"\beta"]\arrow[dr,"\gamma"'] \& {}\\
{} \& {} \& c
\end{tikzcd}}
dans une catégorie $\deleted{\unsure{X}}Z\in\mathrm{Ob}\,\mathrm{Cat}$, dont la limite inductive \inserted{$b\amalg_a c$} dans $Z^{\wedge}$ soit le préfaisceau final $e_Z$, i.e. tel que l’on ait
\[
b(z)\amalg_{a(z)}c(z)\overset{\sim}{\longrightarrow}\text{ ensemble ponctuel}
\]
pour tout objet $z$ de $Z$ (où $a(z)=\mathrm{Hom}_{Z}(\cdot,z)$), sans que pour autant $Z$ soit asphérique, ni même $1$-connexe\textsuperscript{*}.
% slanted
\margin{\textsuperscript{*}\uline{faux} -- il est $1$-connexe forcément, mais non $2$-connexe\unsure{;|.} (cf commentaire \uline{XII} p. \unsure{101})}
L’exemple le plus simple que j’ai trouvé, est l’exemple ``universel'' d’un diagramme (a), où \deleted{les deux} ni l’une ni l’autre des deux flèches $a\unsure{\text{---}|\longrightarrow}b$, $a\text{---}c$ n’est mono (car si elle l’est, alors on peut prouver que $Z$ est bel et bien asphérique, comme nous verrons par la suite). C’est la catégorie
\diagram{\begin{tikzcd}[ampersand replacement=\&,row sep=small,column sep=small]
{} \& {} \& {} \& b_0\\
Z=\biggl\{ \& \unsure{\sigma_0|o_0}\arrow[r,bend left=16,"\lambda_0"]\arrow[r,bend right=16,"\mu_0"'] \& a_0\arrow[ur,"\beta_0"]\arrow[dr,"\gamma_0"'] \& {} \& \biggr\}\\
{} \& {} \& {} \& c_0
\end{tikzcd}}

% archive sheet number: 345
\scanpage{346}
\mspage{\deleted{\unsure{12|10}}11}

déduite de $\unsure{\Upsilon|Y}$ en lui rajoutant un objet $\unsure{o_0|\sigma_0}$, relié à $a_0$ par deux flèches distinctes $\lambda_0,\mu_0$ telles que
\[
\beta_0\lambda_0=\beta_0\mu_0\qquad\text{et}\qquad\gamma_0\lambda_0=\gamma_0\mu_0,
\]
ces deux flèches composées étant, avec les quatre flèches $\beta_0,\gamma_0,\lambda_0,\mu_0$ marquées sur le diagramme, les seules (six) flèches non identiques de $Z$. $\Upsilon$ \deleted{\unsure{s’identifie|identifiée}} est immergée dans $Z$ \deleted{\unsure{comme un|pour un}} un fermé, de complémentaire la catégorie ponctuelle réduite à $o_0$. (N.B. L’immersion fermée $\Upsilon\longrightarrow Z$ n’est pas immersion de Dwyer \inserted{(\unsure{comme} $Z$ est le plus petit ouvert de $Z$ contenant $\Upsilon$)} i.e. n’est pas coasphérique, car la catégorie $o\backslash\Upsilon$, cofibrée sur $\Upsilon$ avec fibres discrètes, de cardinal $2$ en $a_0$ et $1$ en $b_0$ et $c_0$, est la catégorie
\diagram{\begin{tikzcd}[ampersand replacement=\&,row sep=small,column sep=large]
a'_0\arrow[r]\arrow[dr] \& \widetilde b_0\\
a''_0\arrow[ur]\arrow[r] \& \widetilde c_0
\end{tikzcd}}
, \deleted{\unsure{qui}} \inserted{i.e.}
\diagram{\begin{tikzcd}[ampersand replacement=\&,row sep=tiny,column sep=tiny]
{} \& \bullet\\
\bullet\arrow[ur]\arrow[dr] \& {} \& \bullet\arrow[ul]\arrow[dl]\\
{} \& \bullet
\end{tikzcd}}
\inserted{, qui} a le type d’homotopie du cercle). Ni l’immersion ouverte $\{o_0\}\unsure{\text{---}|\longrightarrow}Z$, car $\{o_0\}/a_0$ est la cat. discrète de card. $2$.

\deleted{\unsure{Z}}L’inclusion \deleted{\unsure{de}} $i:Z_0\text{---}Z$ de
\diagram{\begin{tikzcd}[ampersand replacement=\&,column sep=small]
Z_0={} \& o_0\arrow[r,bend left=20,"\lambda_0"]\arrow[r,bend right=20,"\mu_0"'] \& a_0
\end{tikzcd}}
dans $Z$ \unsure{est} \deleted{\unsure{une} \unsure{équivalence}} homotopisme\textsuperscript{*} (car il existe une rétraction évidente $r$ de $Z$ sur $Z_0$\textsuperscript{*}, et une flèche $i\unsure{r}\longrightarrow\mathrm{id}_Z$), donc $Z$ est \unsure{homotope} à \unsure{$Z$|$Z_0$}, qui a le type d’homotopie d’un cercle, donc n’est pas $1$\inserted{$\unsure{2}$}-connexe,
% slanted
\margin{faux, il n’y a pas de rétraction du tout ! En fait, $Z$ est la \uline{suspension} de
\diagram{\begin{tikzcd}[ampersand replacement=\&,row sep=tiny,column sep=small]
{}\arrow[r,bend left=25]\arrow[r,bend right=25] \& {}
\end{tikzcd}}
, qui \unsure{a le|est de} type d’homotopie \unsure{(topologique)} du cercle, \unsure{et} \unsure{le type} d’homotopie de \unsure{la} $2$-sphère.}

% archive sheet number: 346
\scanpage{347}
\mspage{12\deleted{\unsure{9}}}

D’autre part, la vérification de
\[
b_0\amalg_{a_0}c_0\unsure{=|\simeq}e_Z\qquad\text{dans }Z^{\wedge}
\]
est immédiate.

On va utiliser cette flèche
\[
\unsure{\Upsilon|Y}\longrightarrow Z
\]
pour construire des exemples intéressants (et ``\unsure{déchirants|délirants}) de sommes amalgamées
\diagram{\begin{tikzcd}[ampersand replacement=\&]
U\arrow[r,hook,"i"]\arrow[d,"p"'] \& X\arrow[d,"q"]\\
U'\arrow[r,hook,"i'"] \& X',
\end{tikzcd}}
avec $i$ \deleted{\illegible{1}} imm. ouverte, \deleted{\illegible{1}} telles qu’on ait même
\[
i:U\longrightarrow\overline U\qquad\text{est asphérique, donc }\in W,
\]
\deleted{\unsure{sans}} \inserted{alors} que pourtant on \deleted{ait} aura
\[
i':U'\unsure{\text{---}|\longrightarrow}\overline{U'}\qquad\notin W,
\]
et ce qui mieux, (ou pis) est,
\[
\deleted{\unsure{q}}p\in W\qquad\text{mais}\qquad\deleted{\unsure{p}}q\notin W.
\]
(et éventuellement, en $=$ temps, $i\in W$ mais $i'\notin W$). De plus (contrairement aux constructions des \unsure{prop.|pp.} 10, 11 (p. 101-107), les catégories envisagées \deleted{seront} \inserted{pourront toutes être} finies. Je suis particulièrement intéressé par le cas où $X$ est un ``\unsure{carré}''
\[
X=\Phi_n\times\Phi_m\qquad(n,m\geq 2)
\]
et où $U$ est son ``bord''
\[
U=\dot X\overset{\mathrm{def}}{=}(\dot\Phi_n\times\Phi_m)\cup(\Phi_n\times\dot\Phi_m).
\]
\deleted{Nous \unsure{reviendrons|avions vu}, sur \unsure{ce|le} cas de \unsure{suite|tête}, et \unsure{précisons|préciserons} les}

% archive sheet number: 347
\scanpage{348}
\mspage{\unsure{124|$12_4$}}

\deleted{\unsure{cheminer par des}}

\deleted{\unsure{cheminer|cheminons} par des considérations \inserted{$\pm$} générales.

Je suppose que l’on ait dans $U$ une sous-catégorie $V$ ouverte \deleted{ou fermée (elle sera fermée dans le cas qui nous intéresse)}, isomorphe à $\unsure{\Upsilon|Y}$ ; plus précisément $V$ est \deleted{\unsure{relative à}} \inserted{associée à} un point $t_0\in\dot U$ \deleted{comme} par
\[
\begin{aligned}
\text{\textcircled{a}}\quad\biggl\{
&\deleted{\unsure{V=U_{/t}}}\ \unsure{V}=\unsure{U}_{/t_0}\ \deleted{\unsure{et}}\quad\text{si }t\in\dot U\setminus\{t_0\}\\
&\deleted{\unsure{V_{/t_0}\longrightarrow V}\quad\unsure{V_{\!/t}}\ \text{\unsure{est d’un type}}}\ \unsure{\text{(ordinal }1\text{)}|\text{(élément }1\text{)}}.
\end{aligned}
\]
\unsure{(} le foncteur canonique $U_{/t_0}\unsure{\text{---}|\longrightarrow}U$ étant une immersion ouverte (ce qui est automatique si $\overline U$, voire $X$ tout entier, est une catégorie ordonnée).

Je m’aperçois à l’instant que dans le cas qui m’intéresse particulièrement, \deleted{cette} \inserted{la} catégorie \inserted{$U_{/t}$ elle-$=$,} n’est pas \inserted{\unsure{isom à}} $\Upsilon$, \inserted{comme je croyais l’avoir \unsure{lu|vu},} mais la catégorie un peu plus grande
\[
\deleted{\unsure{\Phi_2}}
\]
\diagram{\begin{tikzcd}[ampersand replacement=\&,column sep=0.9em,row sep=tiny]
\Phi_2=\bigl\{ \& {}\arrow[r] \& {} \& {}\arrow[l]\arrow[r] \& {} \& {}\arrow[l] \& \bigr\},\\
{} \& {} \& {} \& \underbrace{\hspace{3.5em}}_{\Upsilon}
\end{tikzcd}}
dans laquelle $\Upsilon$ est immergée en position centrale stratégique. \deleted{Néanmoins, je vais supposer \unsure{que $V$} \illegible{3}} Le point, \deleted{\unsure{crucial}} est d’utiliser seulement \inserted{la structure intrinsèque de} $V$, c’est \inserted{\textcircled{b} $V$ est asphérique et} \textcircled{b} qu’il existe une \deleted{\unsure{fonction}} \inserted{flèche}
\[
\deleted{\unsure{V}}\qquad p_V^0:V\longrightarrow V'_0
\]
dans Cat, telle que
\[
p_{V!}^0(e_V)=e_{V'_0},\qquad V'_0\ \uline{\text{non asphérique}}.
\]
Soit donc $V'$ déduite de $V'_0$ en lui ajoutant un élément final, \unsure{d’où|donc} $V'_0\hookrightarrow\unsure{V'_0|V'}$ \unsure{(imm.|l’imm.} ouverte),}

% archive sheet number: 348
\scanpage{349}
\mspage{13}

\uline{Proposition} : Soient $n,m\geq 2$, et considérons l’inclusion $U=\partial X\overset{i}{\hookrightarrow}X$ du bord du ``carré'' $X=\Phi_n\times\Phi_m$ dans $X$. C’est une immersion ouverte \inserted{$W$-}asphérique, de façon précise les $U_{/t}$ ($t\in\dot U=\overline U\setminus U$) sont isomorphes à $e$ (cat. ponctuelle) ou à
\diagram{\begin{tikzcd}[ampersand replacement=\&,row sep=tiny,column sep=small]
{} \& {} \& {}\arrow[dl]\\
\unsure{\Phi|\Psi}={} \& {}\\
{} \& {} \& {}\arrow[ul]
\end{tikzcd}}
(catégorie avec élément final), sauf \deleted{\unsure{pour}} les quatre éléments $t_i$ des quatre ``coins'' de $\dot U$, où c’est la catégorie
\diagram{\begin{tikzcd}[ampersand replacement=\&,row sep=tiny,column sep=small]
{} \& {} \& a \& a'\arrow[l]\\
\widetilde{\unsure{\Upsilon|Y}}={} \& o\arrow[ur]\arrow[dr] \& {}\\
{} \& {} \& b \& b'\arrow[l]
\end{tikzcd}}
\diagram{\begin{tikzcd}[ampersand replacement=\&,row sep=small,column sep=small]
o\arrow[r]\arrow[d] \& a\arrow[d] \& a'\arrow[l]\arrow[r]\arrow[d] \& {} \& \cdots\\
b\arrow[r] \& \bullet_{t_1} \& {}\arrow[l] \& {} \& \cdots\\
b'\arrow[u]\arrow[r]\arrow[d] \& {}\arrow[u] \& {}\arrow[u]\arrow[l] \& {} \& \cdots\\
{} \& \cdots \& \cdots \& {} \& \cdots
\end{tikzcd}}
n’ayant \deleted{\unsure{un} élément initial, mais} pas d’él\textsuperscript{t} final. Donc $U\xrightarrow{i}X$ n’est pas imm. ouverte de Dwyer, (mais $U\longrightarrow X\setminus\{t_1,t_2,t_3,t_4\}$ l’est...). Ceci posé, je dis que $i:U\longrightarrow X$ n’est pas une $W$-cofibration.

\uline{\unsure{Dém.}} Soit
\[
F\unsure{\subset|\lhook\joinrel\relbar}U
\]
l’inclusion dans $U=\partial X$ de la catégorie-\unsure{coin|coïn}
\diagram{\begin{tikzcd}[ampersand replacement=\&,row sep=small,column sep=small]
o\arrow[r]\arrow[d] \& a\\
b \& {}
\end{tikzcd}}
(cf figure), \inserted{isom. à $\Upsilon$.} \unsure{(}On voit aussitôt que c’est une imm. fermée de Dwyer (le plus petit ouvert de $U$ contenant $\Upsilon$ est précisément la catégorie $\widetilde\Upsilon=U_{/t_1}$ de tantôt), donc une $W$-cofibration. Si $U\unsure{\text{---}|\longrightarrow}X$ était une $W$-cofibration, il en serait de $=$ du composé $V\longrightarrow X$. On va voir que ce n’est pas le cas.

Soit
\diagram{\begin{tikzcd}[ampersand replacement=\&,row sep=small,column sep=small]
\overline V={} \& o\arrow[r]\arrow[d] \& a\arrow[d]\\
{} \& b\arrow[r] \& t_1
\end{tikzcd}}
l’adhérence de $V$ ds $X$,

% archive sheet number: 349
\scanpage{350}
\mspage{14}

et considérons l’inclusion
\[
\overline V\longrightarrow X,
\]
je dis que c’est une imm. fermée de Dwyer, i.e. que pour tout $s\in W\setminus\overline V$, où $W$ est le plus petit ouvert de $X$ contenant $\overline V$, i.e.
\diagram{\begin{tikzcd}[ampersand replacement=\&,row sep=small,column sep=small]
{} \& \unsure{a|o}\arrow[r]\arrow[d] \& a\arrow[d] \& a'\arrow[l]\arrow[d]\\
W={} \& b\arrow[r] \& t \& a''\arrow[l]\\
{} \& b'\arrow[u]\arrow[r] \& b''\arrow[u] \& c\arrow[u]\arrow[l],
\end{tikzcd}}
la catégorie $s\backslash\overline V$ a un objet initial, (les catégories $s\backslash\overline V$ qu’on obtient, à partir des \deleted{quatre} \inserted{cinq} choix possibles pour $s$, sont $\Delta^1$ (si $s=a'$ ou $b'$) ou $e$ (pour $s=a'',b''$ ou $c$). Donc $\overline V\unsure{\text{---}|\longrightarrow}X$ est une $W$-cofibration.

Ceci posé, on sait qu’il existe, pour l’inclusion (imm. ouverte)
\[
V\hookrightarrow\overline V\qquad\text{i.e.}\qquad \unsure{\Upsilon|Y}\subset\deleted{\unsure{\Upsilon}}\Delta^1\times\Delta^1,
\]
un foncteur $V\text{---}V'$, donnant lieu à un carré cocartésien
\diagram{\begin{tikzcd}[ampersand replacement=\&,column sep=small,row sep=small]
{} \& \deleted{\unsure{V}}\arrow[d] \& V\arrow[r,hook,no head,"i"]\arrow[d,"p"'] \& \overline V\arrow[d,no head,"\overline p"]\\
\text{(a)} \& \deleted{\unsure{V'}} \& V'\arrow[r,hook,"i'"] \& \overline{V'}
\end{tikzcd}}
avec
\[
\begin{aligned}
&i,p\in W\quad\deleted{\text{\unsure{(car $V'$ est asph.)}}}
\inserted{\text{(car $V,V',\overline V$ $W$-asph.)}}\quad\text{mais }\inserted{i',}\overline p\notin W
\quad\deleted{\text{\unsure{(i.e. $\overline{V'}$ pas}}}\\
&\hspace{5em}\text{(i.e. $\overline{V'}$ \uline{pas} $W$-asph.\deleted{\unsure{\text{asph.}}}).}
\end{aligned}
\]
Je rappelle qu’on prend d’abord
\diagram{\begin{tikzcd}[ampersand replacement=\&,row sep=small,column sep=small]
{} \& {} \& {} \& a \& {} \& \alpha\lambda=\alpha\mu\\
V'_0={} \& o'\arrow[r,bend left=20,"\lambda"]\arrow[r,bend right=20,"\mu"'] \& o\arrow[ur,"\alpha"]\arrow[dr,"\beta"'] \& {} \& \text{avec} \& {}\\
{} \& {} \& {} \& b \& {} \& \beta\lambda=\deleted{\unsure{\alpha|\beta}}\beta\mu
\end{tikzcd}}
et l’imm. fermée ``identique'' de $V$ dans $V'$
\[
p_0:V\longrightarrow V'_0
\]

% archive sheet number: 350
\scanpage{351}
\mspage{15}

(qui est telle que
\[
p_0^{\deleted{\unsure{*}}}{}_{!}(e_V)=e_{V'_0}\quad)
\]
puis on prend pour $V'$ la catégorie obtenue en \unsure{ajoutant|rajoutant} un objet final à $V'_0$
\diagram{\begin{tikzcd}[ampersand replacement=\&,row sep=small,column sep=small]
{} \& {} \& {} \& a\arrow[dr,"\rho"] \& {} \& {} \& {} \\
V'={} \& o'\arrow[r,bend left=18,"\lambda"]\arrow[r,bend right=18,"\mu"'] \& \unsure{o|a}\arrow[ur,"\alpha"]\arrow[dr,"\beta"'] \& {} \& e \& \text{avec} \& {\left\{\begin{gathered}\alpha\lambda=\alpha\mu\\\beta\lambda=\beta\mu\\\rho\alpha=\sigma\beta\end{gathered}\right.}\\
{} \& {} \& {} \& b\arrow[ur,"\sigma"'] \& {} \& {} \& {}
\end{tikzcd}}
et pour $p$ le foncteur composé
\diagram{\begin{tikzcd}[ampersand replacement=\&,column sep=small]
V\arrow[r,"p_0"]\arrow[rr,bend right=20,"p"',dash] \& V'_0\arrow[r,hook,"{\deleted{\unsure{i}}p_1}"] \& V',
\end{tikzcd}}
où $p_1$ est l’inclusion (ouverte). On trouve \unsure{alors le|dans le} carré (a), \deleted{où $p\in W$ (car $V$ et $V'$ sont $W$-asph.\unsure{)}}
\deleted{\diagram{\begin{tikzcd}[ampersand replacement=\&,row sep=small,column sep=small]
V\arrow[r,hook]\arrow[d,"p"'] \& \overline V\\
V'\arrow[r,hook,no head] \& {}
\end{tikzcd}}}
\deleted{\unsure{vérifiant}, mais $\overline p\notin W$ i.e. (comme $\overline V$ est $W$-asph.) $\overline{V'}$ \uline{non} $W$-asph.,} \deleted{\unsure{avec}} $\overline{V'}$ non $W$-asph, ou ce qui revient au $=$, $i'\notin W$\deleted{\unsure{asph.}}. Car comme $\overline{V'}\setminus V'\simeq\overline V\setminus V$ est la catégorie ponctuelle, $i'\in W$ équivaut à : $V'/t$ $W$-asph. (cf \unsure{cor.} 9, p. 149 dans (\unsure{XII|XIb})),
% slanted
\margin{moyennant ($W$ \unsure{7} bis)}
Or par construction $V'/t\simeq V'_0$, qui n’est pas $W$-asphérique (dans le cas, disons, où $W=\unsure{W_1|W_\infty}$, car $\pi_1(\unsure{\cdot}V'_0)\unsure{\simeq|=}\mathbb Z$).

\deleted{Cela prouve la proposition. \deleted{\unsure{On trouve un}} \inserted{Pour expliciter} ce contre-exemple, notons que l’on aura un}

% archive sheet number: 351
\scanpage{352}
\mspage{16}

\deleted{\unsure{Consi}}Considérons \unsure{donc|alors} le diagramme de carrés cocartésiens
\diagram{\begin{tikzcd}[ampersand replacement=\&,row sep=small,column sep=small]
\deleted{V} \& {} \& {}\\
V\arrow[r,hook,"i",no head]\arrow[d,"p"'] \& \overline V\arrow[r,hook]\arrow[d,"\overline p"] \& X\arrow[d,"q"]\\
V'\arrow[r,hook,"i'",no head] \& \overline{V'}\arrow[r,hook,no head] \& X'
\end{tikzcd}}
\deleted{Comme} $\overline V\longrightarrow X$ est imm. \inserted{fermée} de Dwyer, \deleted{\unsure{de}} le carré de droite est $W$-cocartésien donc \deleted{\unsure{on}}
\[
(\overline p\in W)\Longleftrightarrow(q\in W),
\]
donc \inserted{(comme $\overline p\notin W$)} on a \deleted{\unsure{q}}
\[
p\in W,\quad\text{mais }q\notin W,
\]
ce qui montre bien que $V\unsure{\text{---}|\longrightarrow}X$ n’est pas une $W$-cofibration. Factorisant $V\text{---}X$ par $U$ (au lieu de $\overline V$) on trouve
\diagram{\begin{tikzcd}[ampersand replacement=\&,row sep=small,column sep=large]
V\arrow[r]\arrow[d,"{p\in W}"'] \& U\arrow[r,dash]\arrow[d,"{q_U}"',"{\unsure{U|W}}","{\in}"{description,fill=none}] \& X\arrow[d,"{q\deleted{\unsure{\in}}\notin W}"]\\
V'\arrow[r,dash] \& U'\arrow[r,dash] \& X',
\end{tikzcd}}
où $q_U\in W$ puisque $p\in W$ et $V\longrightarrow U$ est imm. de Dwyer. Cela rend explicite que $U\longrightarrow X$ n’est pas une $W$-cofibration, puisque
\[
q_U\in W\quad\text{mais}\quad q\notin W,
\]
\hfill \unsure{qed.}

% archive sheet number: 352
\scanpage{353}
\mspage{\textcircled{17}}

\uline{Recollement dans Cat d’un ouvert et d’un fermé complémentaire}

\diagram{\begin{tikzcd}[ampersand replacement=\&,column sep=small,row sep=tiny]
X_0\arrow[r,hook,no head] \& X \& X_1\arrow[l,hook,no head]\\
\text{ouv.} \& {} \& \text{fermé compl.}
\end{tikzcd}}
dans Cat
\[
\widetilde X_{01}\unsure{\subset|\in}\mathrm{Fl}(X)\qquad
\text{\parbox{.48\linewidth}{sous-cat. pleine des flèches de source dans $X_0$, but dans $X_1$}}
\]
\diagram{\begin{tikzcd}[ampersand replacement=\&,row sep=large,column sep=large]
{} \& {} \& \widetilde X_{01}\arrow[dl,"{\substack{p\\\deleted{\text{\unsure{co}}}\text{fibrant}\\\text{fibres les }x_0\backslash X_1}}"']\arrow[dr,"{\substack{q\\\text{cofibrant, fibres les }X_{0/x_1}}}"]\\
\text{(*)} \& X_0\arrow[dr,"i"',""{name=L,inner sep=0pt}] \& {} \& X_1\arrow[dl,"j",""{name=R,inner sep=0pt}]\\
{} \& {} \& X
\arrow[from=L,to=R,bend left=16,shorten <=4pt,shorten >=4pt]
\end{tikzcd}}

\uline{Question} Ce diagramme est-il $W$-cocartésien ?

OK si $i$ \inserted{$W$-}asp. ou $j$ \inserted{$W$-}coasphérique, i.e. $X\longrightarrow\Delta^1$ lisse ou propre -- c’est un cas où on a soit $X_0\xrightarrow{W}X$, soit $X_1\xrightarrow{W}X$. Mais $=$ en supposant $i$ ou $j\in W$, et de plus $X$ asphérique (donc $X_0$ et $X$ asph, $X_1$ et $X$ asph), et $W=\mathbb W_\infty$, je n’ai pu prouver que le diagramme est cocartésien \inserted{i.e. $\unsure{q\in|\int}W$ \unsure{resp.|imp.|asp.} $p\in W$\textsuperscript{*}}.
% slanted
\margin{\textsuperscript{*} Je suppose que $W$ satisfait Loc(6)}
(Pas non plus si $X_0\simeq e$, $X,X_1$ \inserted{$W$-}asph $\Longrightarrow {}_e\backslash X$ \inserted{$W$-}asph i.e. $X_1\unsure{\text{---}|\longrightarrow}X$ $W$-coasphérique, ..) Mais j’ai présumé du moins, \unsure{dans|ds} tous les cas, que
\diagram{\begin{tikzcd}[ampersand replacement=\&,row sep=tiny,column sep=small]
{} \& {} \& {} \& X_0\\
\text{(**)} \& \unsure{\mathcal X|\mathscr X}\overset{\mathrm{def}}{=}\displaystyle\int \& \widetilde X_{01}\arrow[ur]\arrow[dr] \& {}\arrow[r,"\varphi"] \& X\\
{} \& {} \& {} \& X_1
\end{tikzcd}}
\deleted{induit une \unsure{trans.}} est dans $\mathbb W_1$. \unsure{$\Big]$} Chose \uuline{fausse} \inserted{(sauf si $i$ ou $j\in W=\mathbb W_\infty$,)} si on remplace $\widetilde X_{01}$ par $\widetilde X'_{01}$, formé des triples $(x_0,x_1,c)$, avec $x_0\in\mathrm{Ob}\,X_0$, $x_1\in\mathrm{Ob}\,X_1$, $c:x_0\dashrightarrow x_1$ un \uline{chemin} de $x_0$ à $x_1$ dans $X$ ! \unsure{$\Big\}$} On a

% archive sheet number: 353
\scanpage{354}

% printed form
% letterhead
% Four long pen strokes cancel the printed heading and addressee fields.
\noindent\begin{minipage}[t]{.48\linewidth}\centering
\deleted{MINISTERE DE L’EDUCATION NATIONALE}\\
\deleted{ACADEMIE D’AIX-MARSEILLE}\\
\rule{2em}{.3pt}\\[.5em]
\deleted{Inspection Académique}\\
\deleted{de Vaucluse}\\
\rule{2em}{.3pt}\\[.5em]
\deleted{Tél. : 90 82 99 76}\\[.5em]
\deleted{Référence : Secrétariat n°}
\end{minipage}\hfill
\begin{minipage}[t]{.48\linewidth}
\deleted{AVIGNON, le}\\[2.5em]
\deleted{L’Inspecteur d’académie,}\\
\deleted{directeur des services départementaux}\\
\deleted{de l’Education Nationale de Vaucluse,}\\[2em]
\deleted{à M}
\end{minipage}

% archive sheet number: 354
\scanpage{355}
\mspage{\textcircled{18}}

donc une situation ``topologique'' tout à fait spéciale à la catégorie des modèles Cat, qui ne semble pas avoir d'analogue dans des catégories telles que Esp ou $\Delta^{\wedge}$, où pourtant les notions d'ouvert et de fermé complémentaire (avec grains de sel \unsure{ds|dans} $\Delta^{\wedge}$) ont un sens.

L'argument \inserted{catégorique} quasi-trivial qui établit que $(**)$ \deleted{induit} est dans $\mathbb W_1$, i.e. induit une équivalence de catégories entre $\mathrm{Rev}(X)$ et $\mathrm{R\acute{e}v}(\unsure{\mathcal X|\mathscr X})\overset{\sim}{\simeq}\mathrm{Rev}(X_0)\mathop{\times}\limits^{2}_{\mathrm{Rev}(\widetilde X_{01})}\mathrm{Rev}(X_1)$, \unsure{me|ne} semble devoir s'étendre aux dimensions supérieures, à coups de $n$-gerbes \inserted{sur $X$, \unsure{..}} -- mais les fondements sur la théorie des $n$-gerbes \deleted{\unsure{de}} manquent, (Et je n'ai pas le courage d'expliciter l'argument pour $n=2$ (où pourtant Giraud fournit ce qu'il faut.) Aussi je conjecture que $\varphi$ est \deleted{tout} au moins dans $\mathbb W_\infty$, i.e. que le diagramme $(*)$ est $W$-cocartésien si $W\unsure{>|\supset}\mathbb W_\infty$ (et satisfait $\mathrm{Loc}(6'\unsure{.})$).
% slanted
\margin{Mais j'ai fait \unsure{quand même} une esquisse de vérification pour $n=2$, qui me convainc \unsure{qu'on a} bien $\varphi\in\mathbb W_2$, \unsure{Et je|et} \unsure{m’en suis|j’ai} acquis \unsure{de la} \unsure{sorte|certitude} qu'on a même $\varphi\in\mathbb W_\infty$.}

% archive sheet number: 355
\scanpage{356}

% printed form
% letterhead
\noindent\begin{minipage}[t]{.48\linewidth}\centering
MINISTERE DE L’EDUCATION NATIONALE\\
ACADEMIE D’AIX-MARSEILLE\\
\rule{2em}{.3pt}\\[.5em]
Inspection Académique\\
de Vaucluse\\
\rule{2em}{.3pt}\\[.5em]
Tél. : 90 82 99 76\\[.5em]
Référence : Secrétariat n°
\end{minipage}\hfill
\begin{minipage}[t]{.48\linewidth}
AVIGNON, le\\[2.5em]
L’Inspecteur d’académie,\\
directeur des services départementaux\\
de l’Education Nationale de Vaucluse,\\[2em]
à M
\end{minipage}

% archive sheet number: 356
